\chapter{Conclusion générale et perspectives}
\label{chap:conclusion}

\section{Synthèse du travail réalisé}

Ce projet a permis de concevoir et développer \projectname, une plateforme
qui étend une installation Wazuh existante d'une couche d'investigation
automatisée fondée sur un agent d'intelligence artificielle. Le travail
réalisé couvre l'ensemble de la chaîne~: collecte périodique des alertes,
moteur de politiques configurable, boucle agentique de raisonnement outillé
inspirée du paradigme ReAct \cite{react2022}, traçabilité complète des
actions de l'IA, assistant conversationnel partageant les mêmes outils,
enrichissement par une source de renseignement externe (réputation d'IP),
et authentification déléguée aux comptes Wazuh existants. L'ensemble a été
validé fonctionnellement contre des services réels tout au long du
développement (\cref{chap:resultats}).

\section{Difficultés rencontrées}

Plusieurs difficultés ont jalonné le développement, et ont chacune donné
lieu à une évolution de conception~:

\begin{itemize}
  \item les quotas de débit (\textit{tokens per minute}) imposés par le
        fournisseur de LLM utilisé en développement se sont révélés plus
        contraignants que prévu, obligeant à réduire la taille du contexte
        transmis à chaque appel (invite système, contexte de référence,
        résultats d'outils) sans dégrader la qualité du raisonnement~;
  \item certains modèles ont ponctuellement tenté d'appeler un outil alors
        qu'il leur était explicitement demandé de conclure
        (\texttt{tool\_choice="none"}), ce qui a nécessité la détection et
        la gestion explicite de ce cas plutôt que de le traiter comme une
        erreur fatale~;
  \item les arguments transmis par le modèle lors d'un appel d'outil ne
        respectent pas toujours strictement le type attendu (un nombre
        transmis sous forme de texte), ce qui a conduit à une normalisation
        systématique des arguments avant exécution (\cref{lst:coerce})~;
  \item la conception d'un mécanisme d'authentification sans dupliquer la
        base de comptes existante a demandé une réflexion spécifique,
        résolue par la délégation de la vérification au Wazuh Indexer.
\end{itemize}

Ces difficultés, bien que ponctuellement contraignantes, ont eu un effet
positif sur la robustesse finale du système~: chacune a été traitée comme
un cas à absorber structurellement plutôt que comme une exception isolée.

\section{Perspectives}

Plusieurs axes d'évolution ont été identifiés pour la suite du projet~:

\begin{itemize}
  \item \textbf{Gestion fine des permissions}~: exploiter le champ de rôle
        déjà présent sur les utilisateurs pour restreindre certaines
        actions (par exemple la gestion des politiques) aux seuls profils
        autorisés~;
  \item \textbf{Migration vers un modèle de langage auto-hébergé}~:
        l'abstraction \texttt{LLMProvider} a précisément été conçue pour
        permettre cette transition sans modification de la logique
        métier, afin de s'affranchir des contraintes de quota et de coût
        d'une API distante~;
  \item \textbf{Élargissement des sources de données}~: intégration
        d'autres capteurs (Suricata pour la détection réseau, journaux
        \textit{cloud} de type CloudTrail) sous forme de nouveaux outils,
        sans remise en cause de l'architecture~;
  \item \textbf{Enrichissement par RAG}~: adossement de l'agent à une base
        de connaissances de menaces (par exemple structurée autour du
        référentiel MITRE ATT\&CK \cite{mitre_attack}) pour améliorer la
        qualification des verdicts~;
  \item \textbf{Suite de tests automatisés}~: couverture unitaire des
        outils et de la boucle agentique, et tests d'intégration de l'API,
        pour sécuriser les évolutions futures~;
  \item \textbf{Possibilité d'annuler une investigation en cours}, et
        \textbf{tableaux de bord analytiques} (tendances de charge
        d'alertes, taux de faux positifs observé dans le temps).
\end{itemize}

En définitive, ce projet démontre la faisabilité d'une automatisation du
triage d'alertes de sécurité par un agent fondé sur un modèle de langage,
tout en conservant l'exigence de traçabilité et de sobriété d'accès aux
données propre à un contexte de cybersécurité — les deux axes qui ont guidé
l'ensemble des choix de conception présentés dans ce rapport.
